%% LyX 2.4.2.1 created this file.  For more info, see https://www.lyx.org/.
%% Do not edit unless you really know what you are doing.
\documentclass[american]{article}
\usepackage[T1]{fontenc}
\usepackage[utf8]{inputenc}
\usepackage[skip=\medskipamount]{parskip}
\usepackage{amsmath}
\usepackage{amssymb}
\usepackage[a4paper]{geometry}
\geometry{verbose,tmargin=2cm,bmargin=2cm,lmargin=2cm,rmargin=2cm}
\usepackage{babel}
\begin{document}
\title{A short note on Variance Swaps}
\author{Quasar Chunawala}
\date{07-May-2026}
\maketitle
\begin{abstract}
A short note on Variance Swaps and volatility derivatives
\end{abstract}

\section*{Replicating a twice-differentiable payoff with a static position
in vanilla options}

Any twice differentiable payoff $g(S_{T})$ at time $T$ may be replicated
using a static portfolio of European options expiring at time $T$.

From the Breeden-Litzenberger formula, we know that we may write the
\textit{pdf }of the stock-price $S_{T}$ as:

\begin{equation}
p(S_{T},T;S_{t},t)=\left.\frac{\partial^{2}\tilde{C}(K,T,S_{t},t)}{\partial K^{2}}\right|_{K=S_{T}}=\left.\frac{\partial^{2}\tilde{P}(K,T,S_{t},t)}{\partial K^{2}}\right|_{K=S_{T}}
\end{equation}

where $\tilde{C}$ and $\tilde{P}$ represent the undiscounted call
and put prices respectively. 

Then, the value of a claim with a generalized payoff $g(S_{T})$ at
time $T$ is given by:

\begin{align*}
\mathbb{E}\left[g(S_{T})|S_{t}\right] & =\int_{0}^{\infty}\underbrace{p(K,T|S_{t},t)}_{\text{Transition probability}}\times g(K)dK
\end{align*}

For any payoff $g(S_{T})$, the sifting property of a Dirac-delta
function implies that:

\begin{align*}
Ig(S_{T}) & =\int_{0}^{\infty}g(K)\delta(S_{T}-K)dK\\
 & =\int_{0}^{F}g(K)\delta(S_{T}-K)dK+\int_{F}^{\infty}g(K)\delta(S_{T}-K)dK
\end{align*}

In particular, the dirac-delta function has symmetry. $\delta(x)=\delta(-x)$.
The Heaviside step function $\theta(x)=\int_{-\infty}^{x}\delta(y)dy$.
The Heaviside function satisfies $\theta(x)=-\theta(-x)$. Thus,

\begin{align*}
\delta(K-S_{T}) & =\frac{\partial}{\partial K}\theta(K-S_{T})=-\frac{\partial}{\partial K}\theta(S_{T}-K)=
\end{align*}

We integrate each integrand by parts. The first integral simplifies
to:

\begin{align*}
\int_{0}^{F}g(K)\delta(S_{T}-K)dK & =\left[g(K)\left(\int\delta(S_{T}-K)dK\right)\right]_{K=0}^{K=F}-\int_{0}^{F}g'(K)\left(\int\delta(S_{T}-K)dK\right)dK\\
\\\\\\
\end{align*}

The second integral simplifies to:

\begin{align*}
\int_{F}^{\infty}g(K)\delta(K-S_{T})dK & =\left[g(K)\int\delta(S_{T}-K)dK\right]_{K=F}^{K=\infty}-\int_{0}^{F}\left(\int\delta(S_{T}-K)dK\right)g'(K)dK\\
\\
\end{align*}

\end{document}
